\section{Background}
\label{sec:background}
 
\subsection{Paired Program Transformations}
\label{sec:bg-pairs}
 
Many software engineering tasks come in pairs.
Translation has back-translation.
Compilation has decompilation.
Obfuscation has deobfuscation.
Research on bidirectional transformations studies such pairs in a formal way~\cite{czarnecki2009bx}.
Lenses are one well-known formalism in this line~\cite{foster2007lenses}.
A lens couples a forward function with a backward function.
Round-trip laws tie the two functions together.
These laws ask for exact recovery of the source.
 
Learned code transformations rarely meet this condition.
Identifier renaming discards the original names.
Deobfuscation tools can only predict plausible names~\cite{raychev2015jsnice,bichsel2016deguard}.
Many programs are therefore valid deobfuscations of one obfuscated program.
Some task pairs are also not inverses in any strict sense.
Mutation testing injects faults on purpose~\cite{jia2011mutation}.
A mutant thus differs in behavior from its source.
Its repair does not undo a semantics-preserving step.
 
Prior work gives no definition of an inverse pair for learned code transformations.
We use the following one.
Let $\mathcal{P}$ be a set of programs.
Let $p \equiv q$ denote that $p$ and $q$ pass the same test suite.
Let $d$ be a distance over program form, such as lexical or structural similarity.
 
\begin{definition}[Forward/inverse pair]
\label{def:pair}
A transformation $T:\mathcal{P}\rightarrow\mathcal{P}$ is a valid forward task if $T(p) \equiv p$ for every $p \in \mathcal{P}$.
The inverse task of $T$ takes $T(p)$ as input.
A solution is any program $q$ with $q \equiv T(p)$ and $d(q,p) < d(T(p),p)$.
\end{definition}
 
The definition keeps semantic equivalence.
It replaces exact recovery with measurable progress toward the source form.
It accepts many correct inverse outputs.
Obfuscation, translation, and compilation satisfy it.
Mutation and repair do not.
A mutant changes behavior.
BRF (Section~\ref{sec:theorical}) applies to the pairs that satisfy Definition~\ref{def:pair}.
 
\subsection{Obfuscation as a Controlled Testbed}
\label{sec:bg-obfuscation}
 
Obfuscation turns a program into an equivalent program that is harder to understand~\cite{collberg2002watermarking}.
Collberg and Thomborson classify obfuscating transformations by the program aspect they change.
Obfuscation is common in practice.
Wermke et al. measure its use across a large set of Google Play apps~\cite{wermke2018obfuscation}.
Dong et al. study obfuscation techniques in Android apps in the wild~\cite{dong2018android}.
Both studies report identifier renaming as the most common technique.
Dong et al. also observe string encryption in real apps.
Our three obfuscation types follow these observations (Section~\ref{sec:framework_instantiation}).
 
Deobfuscation has a long research history.
Yadegari et al. simplify obfuscated executables through semantics-based analysis~\cite{yadegari2015generic}.
JSNice predicts names and types for minified JavaScript~\cite{raychev2015jsnice}.
DeGuard recovers identifier names in obfuscated Android apps~\cite{bichsel2016deguard}.
DIRE predicts variable names in decompiled code~\cite{lacomis2019dire}.
Schrittwieser et al. survey how well obfuscation resists code analysis~\cite{schrittwieser2016protecting}.
 
This history explains our choice of testbed.
The forward direction is cheap to supply with existing obfuscators.
Forward training pairs are therefore abundant.
The inverse direction must recover lost information.
Reverse supervision is therefore scarce~\cite{10.1145/3728886}.
The three types range from lexical to semantic changes.
This range lets us relate the size of the effect to the difficulty of the inverse.
The translation and compilation checks extend the study to two other pairs (Section~\ref{sec:transl-compil}).
 
\subsection{Adapting Code LLMs to One Task}
\label{sec:bg-adaptation}
 
LLMs now support a wide range of software engineering tasks~\cite{hou2024llm4se}.
Practitioners often adapt a pre-trained model to a single task.
Parameter-efficient fine-tuning makes this adaptation affordable~\cite{hu2022lora}.
LoRA trains low-rank updates.
It keeps the base weights frozen. 
LoRA also forgets less of the base model than full fine-tuning~\cite{biderman2024lora}.
Weyssow et al. compare such techniques for code generation~\cite{weyssow2025peft}.
 
Adapted models are then shared and reused.
Jiang et al. study pre-trained model reuse on Hugging Face~\cite{jiang2023reuse}.
They report missing model attributes.
They also report gaps between claimed and actual performance.
A reused checkpoint thus comes with limited information about what it can still do.
%Evaluation after adaptation usually reports the target task.
Adaptation studies typically report the target task, for example code generation~\cite{weyssow2025peft} or decompilation~\cite{tan2024llm4decompile}.
This paper asks what one-direction adaptation does to the inverse task.
 
\subsection{Motivation}
\label{sec:bg-motivation}
 
% Prior work offers three relevant findings.
% Autoregressive models can learn a relation in one direction~\cite{berglund2024reversal,lv2024reversal}.
% Fine-tuning can erase or hide earlier abilities~\cite{mccloskey1989catastrophic,kirkpatrick2017overcoming,kotha2024implicit}.
% Code models are also affected by forgetting~\cite{weyssow2023continual}.
% Both directions of a task can serve as an evaluation signal~\cite{allamanis2024roundtrip,gu2024cruxeval}.

Prior work offers four relevant findings.
First, autoregressive models can learn a relation in one direction~\cite{berglund2024reversal, lv2024reversal, wang2025binding}.
Second, fine-tuning can erase or hide earlier abilities~\cite{mccloskey1989catastrophic, kirkpatrick2017overcoming, kotha2024implicit}.
Code models are also affected by such forgetting~\cite{weyssow2023continual}.
Third, fine-tuning can cause side effects beyond the target task.
Benign fine-tuning weakens safety alignment~\cite{qi2024finetuning}.
Narrow fine-tuning can cause broad misalignment~\cite{betley2025emergent}.
Fourth, both directions of a task can serve as an evaluation signal~\cite{allamanis2024roundtrip, gu2024cruxeval}.

 
These findings come from separate settings.
The reversal curse is studied on factual recall.
Forgetting is studied on tasks unrelated to the new one.
Round-trip evaluation is applied to a model that has not been adapted.
Studies of LLMs on translation, decompilation, and obfuscation evaluate the direction they target~\cite{pan2024lost,tan2024llm4decompile,mohseni2025obfuscate}.
To our knowledge, no study measures the inverse of a code transformation after fine-tuning on its forward direction.
 
This question has practical weight.
Adapted checkpoints are shared and reused with incomplete documentation~\cite{jiang2023reuse}.
Security analysis relies on both obfuscation and deobfuscation~\cite{wermke2018obfuscation,bichsel2016deguard}.
Execution tests alone cannot reveal a failed inverse.
A copy of the obfuscated input is equivalent to the source (Definition~\ref{def:pair}).
It therefore passes the same tests.
 
% Our study complements these lines of work.
% It combines existing ideas in one protocol.
% It borrows the two-direction view from round-trip evaluation.
% It borrows the held-out control from forgetting studies.
% It adapts reverse training and contrastive learning as recovery probes~\cite{golovneva2024reverse,jain2021contracode}.
% The study covers one controlled testbed and two additional paired settings.
% Its evidence on mechanism is behavioral.
% Within this scope, it asks three questions.
% Does forward fine-tuning preserve the inverse (RQ1)?
% Is the loss specific to the inverse (RQ2)?
% Can explicit supervision or contrastive training recover it (RQ3)?

 Our study addresses this gap.
We fine-tune code LLMs on one direction of a paired transformation.
We then measure both directions, before and after fine-tuning.
The difference between the two measurements reveals what adaptation removed.
We call this protocol the Bidirectional Reasoning Framework (BRF).
% It acts as a regression test for model adaptation.
% Unrelated tasks separate a loss that targets the inverse from general forgetting.
% Recovery strategies test whether the lost direction can be restored.
% Additional paired settings test whether the pattern extends beyond the main testbed.
%Our evidence on mechanism is behavioral.
Within this scope, we ask three questions.
Does forward fine-tuning preserve the inverse (RQ1)?
Is the loss specific to the inverse (RQ2)?
Can explicit supervision or contrastive training recover it (RQ3)?

Three lines of research are closest to ours.
The reversal curse describes a directional asymmetry in factual recall~\cite{berglund2024reversal}.
We study this asymmetry on executable transformations that admit many valid outputs, and we observe it after downstream fine-tuning.
Round-trip correctness also runs a model in both directions~\cite{allamanis2024roundtrip}.
It evaluates a fixed model, whereas BRF compares the two directions before and after adaptation.
Forgetting studies measure what fine-tuning removes from unrelated tasks~\cite{kotha2024implicit, weyssow2023continual}.
We measure what it removes from the inverse of the trained task.